\documentclass[sigconf]{acmart}

\usepackage{amsmath}

\AtBeginDocument{%
  \providecommand\BibTeX{{Bib\TeX}}}

\setcopyright{none}
\settopmatter{printacmref=false}

\begin{document}

\title[Revisiting Model Capacity in Data-Constrained Music Retrieval]{Limited and Task-Dependent Gains in Candidate Ranking Quality}

\author{Keita Katsumi}
\affiliation{%
  \institution{California State Polytechnic University, Pomona}
  \city{Pomona}
  \state{California}
  \country{USA}
}
\email{kkatsumi@cpp.edu}

\renewcommand{\shortauthors}{Katsumi}

\begin{abstract}
This paper investigates the effect of embedding-model capacity on candidate ranking quality in narrow-domain, data-constrained music retrieval settings. While large pretrained audio models have shown strong performance in benchmark tasks, their computational cost raises an important question: whether larger models provide meaningful gains for retrieval-based systems under limited data conditions. We conduct a controlled evaluation of CNN-based and Transformer-based audio representations across capacity tiers on a small-scale classical music archive. Retrieval is formulated as an offline ranking task using metadata-derived proxy relevance labels, and performance is measured using standard information retrieval metrics at multiple cutoff levels. Our results show that, within this constrained setting, increasing model capacity does not consistently improve ranking effectiveness. While capacity-related differences are observed for structured retrieval tasks, these improvements are limited in magnitude, saturate quickly, and remain consistent across ranking depths (K=5 and K=10). In contrast, no statistically significant differences are observed for more abstract, automatically-derived proxy tasks. These findings indicate that, in small-scale and domain-constrained retrieval scenarios, ranking performance is influenced more by the structure and reliability of proxy relevance labels than by model capacity alone. From a system perspective, increased model capacity introduces substantial computational overhead without consistent performance gains, suggesting that smaller or mid-sized models can provide a more favorable cost-performance trade-off. Overall, this study highlights the importance of aligning model capacity with task structure, relevance definition, and dataset characteristics when designing efficient and scalable retrieval systems.
\end{abstract}
\maketitle
\pagestyle{plain}

% Section 1: Introduction
\section{Introduction}\label{sec:introduction}

Modern music retrieval and recommendation systems are typically developed and evaluated using large-scale datasets with abundant user interaction signals. 
However, many real-world applications operate in low-resource environments, such as curated music archives, where both dataset size and user feedback are inherently limited. 
Classical music collections such as the iPalpiti archive exemplify this setting: they contain high-quality, stylistically coherent recordings but lack dense interaction data, requiring content-based approaches to support search and discovery.

In such scenarios, audio embeddings derived from pretrained models provide a practical basis for similarity-based retrieval. 
Recent advances in pretrained audio representations, including convolutional and transformer-based architectures, have demonstrated strong performance on large-scale benchmark datasets. 
However, it remains unclear whether the benefits of increased model capacity transfer to small, domain-constrained collections. In particular, limited data diversity and restricted candidate sets may reduce the effectiveness of higher-capacity models, raising an important design question: to what extent does increasing embedding-model capacity improve ranking quality under low-resource conditions?

Under these constraints, the problem can be viewed as a multimedia retrieval task focused on efficient candidate generation, where both representational quality and computational cost are relevant considerations. In practical recommender systems, content-based embeddings are often used to support candidate generation in cold-start scenarios, making their behavior in data-constrained settings an important but underexplored issue.

This work investigates this question by focusing on the backend embedding and retrieval stage, independent of end-to-end user modeling. Rather than assuming that larger models necessarily yield better performance, we adopt a hypothesis-driven approach to examine whether increasing embedding-model capacity leads to consistent improvements in similarity-based retrieval within a small-scale, single-domain dataset. The dataset represents an extreme low-resource setting, consisting of 203 tracks, and serves as a controlled micro-scale testbed for examining retrieval behavior under data-constrained conditions.

To this end, we evaluate ranking performance across multiple capacity tiers within two architectural families (CNN-based and Transformer-based models). 
Retrieval is formulated as an offline ranking task using metadata-derived proxy relevance labels designed to approximate musically meaningful similarity. Two complementary proxy tasks are considered: a structured metadata-matching task based on composer identity, and a musical-character-based task capturing abstract attributes (e.g., \emph{energetic} vs.\ \emph{calm}). 
Ranking quality is evaluated using standard information retrieval metrics, including NDCG@K and Hit Rate@K.

The study is guided by the following research questions:

\noindent\textbf{RQ1:} How does embedding-model capacity influence ranking quality in content-based retrieval under low-resource conditions?

\noindent\textbf{RQ2:} Under what conditions, if any, do higher-capacity models provide measurable gains over smaller architectures in such settings?

\noindent\textbf{Hypothesis:} In small-scale, single-domain datasets, increasing embedding-model capacity does not consistently improve retrieval effectiveness under the evaluated conditions.

This work makes three main contributions. First, it provides an empirical analysis of embedding-model capacity in low-resource, single-domain retrieval settings. Second, it introduces musically grounded proxy retrieval tasks that enable controlled evaluation of ranking behavior across structured and abstract similarity dimensions. Third, it examines the trade-offs between model capacity and computational cost in constrained environments, providing practical insights into model selection for content-based candidate generation.


% Section 2: Related Work
\section{Related Work}\label{sec:related-work}

Content-based recommendation is especially important in music recommender systems when user interaction data are sparse or unavailable, a setting in which collaborative signals are weak and cold-start effects become severe\cite{schedl2018}. In such cases, content-derived representations of music items provide a practical basis for similarity and recommendation \cite{deldjoo2024}. In such settings, item representations derived from content---including the audio signal itself---become a practical basis for similarity search, retrieval, and downstream recommendation, as content-based methods remain effective for new-item recommendation by relying directly on item descriptors extracted from the music itself \cite{deldjoo2024}. Recent work has shown the effectiveness of pretrained audio representations in music information retrieval while also identifying a gap in their systematic evaluation within music recommender systems \cite{tamm2024}. While end-to-end neural networks are common in recommender systems, recent work has specifically highlighted the utility of evaluating frozen pretrained representations. As a result, the intrinsic quality of backend embeddings is often underexplored. Therefore, an offline comparative analysis of pretrained audio representations is necessary to systematically evaluate their effectiveness in recommendation contexts, independent of downstream model complexity and user interaction noise. This work is closely related to retrieval-based recommendation, where item similarity is used to generate ranked candidate sets prior to downstream personalization. In such systems, embedding models play a key role in candidate generation, directly shaping the quality of retrieved items before any user-aware ranking is applied. This study focuses on the representation and retrieval stage of a two-stage recommender system pipeline, isolating the impact of item embeddings within the candidate generation stage, prior to downstream ranking or personalization models. By doing so, it enables a controlled evaluation of embedding quality in terms of ranking effectiveness.

Prior work has explored a broad range of audio representation architectures, including convolutional neural networks (CNNs), recurrent neural networks (RNNs), transformer-based models, and hybrid architectures \cite{zaman2023}. Historically, CNN and hybrid CNN–RNN models have been highly effective in specific supervised music tasks such as music emotion recognition (MER) \cite{jiang2024}. More broadly, convolutional architectures have played a central role in many music information retrieval tasks, while more recent transformer-based approaches leverage self-attention to model longer-range structure \cite{zaman2023}.In this study, however, the experimental design focuses specifically on pure CNN-based and pure transformer-based model families. CNN-based models remain attractive because their inductive biases support efficient learning of local spectral structure, and CNN-based architectures have demonstrated strong performance across a range of music detection and audio classification tasks, particularly when scaled to larger model capacities \cite{kong2020panns}. Hybrid architectures have also been explored, including CNN--RNN models for temporal modeling \cite{lin2024} and CNN--Transformer models for combining local and long-range representations \cite{pourmoazemi2024}, but they are treated here as background rather than as
primary comparison targets.

Transformer-based models provide a contrasting representation regime. For example, AST applies self-attention directly to spectrograms \cite{gong2021ast}. In practical recommender systems, improvements at the candidate generation stage directly influence downstream ranking quality, making embedding quality a critical factor even in the absence of personalization. Although these models can deliver strong downstream performance, their advantages are typically established on large and diverse datasets, leaving open the question of how well such capacity translates into practical gains in recommendation ranking relative to their computational cost, especially in small, domain-constrained classical archives where models must operate under data sparsity and cold-start constraints. Evaluation strategy is equally important, since proper offline evaluation remains a recognized challenge in music recommender systems research \cite{schedl2018}. Although classification benchmarks are commonly used for audio embeddings \cite{gong2021ast, zaman2023}, strong classification accuracy does not necessarily imply well-structured neighborhoods for similarity-based retrieval \cite{tamm2024}. For retrieval-oriented settings, ranking metrics such as NDCG and top-K recall-based measures are more appropriate, as they evaluate the positions of relevant items and whether they appear within the top-ranked results \cite{krichene2020}. This limitation is consistent with prior observations in MIR research, where datasets are often small or domain-specific, posing challenges for generalizable evaluation \cite{urbano2013}. At the same time, these metrics are sensitive to evaluation design: constrained candidate pools can reduce discriminative power, increase ties, and even alter comparative conclusions \cite{krichene2020, canamares2020}; this is especially important in small archival collections. Diversity and long-tail considerations also remain relevant background concerns in music recommendation, particularly in settings where exploration is valued \cite{porcaro2022}.

Broader recommender-system work addresses personalization through user modeling, sequential
interaction modeling, and contextual signals \cite{lin2025, abbattista2024, schedl2021}, but those
concerns are outside the scope of this research, which isolates backend embedding quality from user modeling and large-scale deployment constraints. Overall, prior work has established that larger and more complex models often improve performance in data-rich settings, yet it remains unclear whether increased embedding-model capacity yields meaningful gains in recommendation ranking quality in small, sparse, single-domain music archives, where content-based representations must compensate for limited behavioral data. This question is particularly important for recommender-system practitioners, as it directly impacts the trade-off between model complexity, computational cost, and ranking effectiveness in real-world cold-start scenarios \cite{schedl2018, deldjoo2024, tamm2024}. This research addresses that gap by examining how capacity interacts with ranking-oriented evaluation under constrained archival conditions \cite{tamm2024, krichene2020, canamares2020}.

Recent work in the recommender systems and information retrieval communities has increasingly focused on data-intensive sequential recommendation models, such as BERT4Rec and related architectures \cite{sun2019bert4rec}, which leverage large-scale user interaction logs. While these approaches have demonstrated strong performance in data-rich environments, they are less applicable in small-scale, data-constrained multimedia archives where interaction data is sparse or unavailable. In such settings, content-based retrieval and representation learning remain essential for candidate generation.Furthermore, it is well established in the information retrieval literature that offline evaluation metrics such as NDCG are sensitive to evaluation design, including the treatment of unjudged items and the construction of relevance signals \cite{yilmaz2008}. This highlights the importance of carefully defining relevance in scenarios where explicit user feedback is not available, as in the present study.


% Section 3: Method
\section{Method}\label{sec:method}\label{ch:problem}\label{ch:methodology}
This section formalizes the retrieval setting and experimental framework used to evaluate the effect of embedding-model capacity on similarity-based retrieval performance.

\subsection{Problem Formulation}

Let $D = \{x_1, x_2, \dots, x_N\}$ denote a collection of $N$ classical music recordings drawn from
the iPalpiti archive. Each item $x_i$ represents a complete audio track.

An embedding model is defined as a function:
\[
f_\theta : \mathcal{X} \rightarrow \mathbb{R}^d
\]
where $\theta$ denotes the model parameters and $d$ the embedding dimensionality. The embedding of
item $x_i$ is:
\[
z_i = f_\theta(x_i)
\]

Similarity between two items is computed using a similarity function $s(\cdot,\cdot)$, such as cosine
similarity:
\[
s(z_q, z_i) = \frac{z_q \cdot z_i}{\|z_q\|\|z_i\|}
\]
Higher similarity values indicate greater embedding-space proximity.

Given a query item $x_q \in D$, all items in the collection except $x_q$ are ranked in descending
order of similarity $s(z_q, z_i)$. In this formulation, each audio track serves as a query, and the remaining tracks form the candidate set to be ranked based on embedding similarity. 

Retrieval quality is evaluated with respect to a relevance function $r_q(x_i) \in \{0,1\}$ defined by domain-specific proxy relevance labels constructed from metadata, and ranking quality is measured using standard IR metrics, with NDCG@K as the primary metric for ranking quality and Hit Rate@K (HR@K) as a complementary measure reflecting whether at least one relevant item appears within the top-K results.

The primary factor of interest is embedding-model capacity within each architectural family. Capacity is operationalized through the predefined small, medium, and large model variants, while architectural family (CNN vs. Transformer) is treated as a separate descriptive factor. Because parameter counts and pretraining conditions are not strictly aligned across families, cross-family comparisons are interpreted descriptively rather than as controlled capacity-matched comparisons. Capacity effects are examined within architectural families via small, base, and large variants under a controlled retrieval setting.

Relevance is defined through two metadata-based domain-specific relevance estimation tasks. Let $g(x)$ denote a metadata labeling
function.

\textbf{Composer Identity Proxy.} Relevance is defined by shared composer identity:
\[
r_q(x_i) = 1 \quad \text{if} \quad g_{\text{composer}}(x_q) = g_{\text{composer}}(x_i)
\]

\textbf{Primary Musical Proxy.} Relevance is defined by shared character labels:
\[
r_q(x_i) = 1 \quad \text{if} \quad g_{\text{character}}(x_q) = g_{\text{character}}(x_i)
\]

For experimental tractability, relevance is treated as binary:
\[
r_q(x_i) \in \{0,1\}
\]
with $r_q(x_i)=0$ when labels do not match. This study assumes that such metadata-derived labels
provide a controlled, though coarse, approximation of musical similarity.
This formulation corresponds to a content-based candidate generation and ranking stage in a recommender system, 
where item embeddings define the retrieval space used to produce ranked candidate sets.

\subsection{Retrieval Pipeline}
The backend evaluation pipeline consists of four stages: audio preprocessing, embedding extraction, similarity computation, and ranking. All components other than embedding extraction are held constant to isolate the effect of embedding capacity.

The dataset is the iPalpiti classical music archive, comprising 203 tracks and 24.9 hours of audio. Ground truth is derived from editorial metadata, including composer identity and character labels used in the domain-specific relevance estimation tasks.

Each complete track serves as the retrieval unit, and embeddings are computed directly from full-length audio inputs without segmentation. Cosine similarity is used for ranking across all models. This controlled design ensures that performance differences are attributable primarily to the embedding representations. To isolate the effect of representation quality, the dataset, query set, candidate pool, similarity function, ranking procedure, and evaluation metrics are kept identical across all models. No downstream personalization, collaborative signal, or learned re-ranking component is applied.

\subsection{Model Setup}
Models are selected to represent discrete capacity tiers within two architectural families: CNN-based and Transformer-based models. Within each family, small, medium, and large variants enable controlled scaling comparisons, while cross-family comparisons are interpreted descriptively because architecture, parameter count, and pretraining objectives may co-vary.

CNN-based models are drawn from a unified convolutional lineage in which scaling is achieved primarily through increased depth under comparable training objectives. Specifically, we use Cnn6 (small), Cnn10 (medium), and Cnn14 (large), following the PANNs framework proposed by Kong et al. \cite{kong2020panns}, to represent controlled capacity scaling within the CNN family.

We select MERT as the representative Transformer-based model due to its strong performance on a wide range of audio understanding tasks and its availability as a standardized pretrained baseline. We use MERT-95M (medium) and MERT-330M (large) to represent different capacity levels within the transformer family, following the MERT framework \cite{li2023mert}. We note that, unlike the CNN family, a small-capacity variant is not included for the Transformer models. This is because publicly available pretrained checkpoints in the MERT framework are provided solely at medium and large scales; as a result, our Transformer-based evaluation is restricted to these available variants.

Note that parameter counts are not strictly aligned across architectural families; therefore, cross-family comparisons are interpreted descriptively rather than as direct capacity-matched evaluations. All selected models provide pretrained checkpoints and support standardized embedding extraction.


\subsection{Proxy Relevance Label Construction}

In practical information retrieval systems, relevance is typically defined based on user interaction logs. However, such signals are unavailable in our setting. To address this limitation, we construct domain-specific proxy relevance labels to approximate ground truth for an item-to-item similarity retrieval task.

Relevance labels for the domain-specific estimation tasks are derived from metadata and a semi-automatic annotation pipeline.

For the Composer Identity Proxy, composer labels are normalized to a canonical form to ensure consistent matching across recordings. This normalization is performed during pseudo-label generation and stored as a stable evaluation field.

For the primary musical proxy, character labels are constructed using a semi-automatic pipeline. Initial label candidates are defined based on the arousal--valence framework \cite{eerola2011}, which represents emotion in a continuous dimensional space. To support this process, pretrained music emotion recognition models are used as auxiliary tools for generating candidate labels, grounded in recent advances in deep learning-based music emotion recognition \cite{jiang2024}.

Ambiguous cases are manually reviewed to ensure consistency across the dataset. The resulting proxy labels are fixed prior to evaluation and reused across all experiments.

Although these proxy labels do not fully capture subjective user preference, they provide a consistent and reproducible evaluation framework for comparing embedding-based retrieval performance in a data-constrained setting.




\subsection{Evaluation Protocol}

Each track in the archive is used as a query in a leave-one-out retrieval setting, where the query item itself is excluded from the candidate pool. Relevance is determined by metadata matching as defined above, and performance metrics are averaged across all query instances.

Because relevance labels are binary and the dataset is relatively small, ranking ties may occur when multiple relevant items share identical similarity scores. To mitigate potential bias from arbitrary ordering, we adopt a tie-aware evaluation strategy by assigning averaged ranks to tied items rather than relying solely on lexicographic ordering.

Performance is primarily assessed using NDCG@5 and Hit Rate@5 (HR@5), with interpretation focused on consistent scaling patterns rather than small numerical differences.

Furthermore, small numerical differences in mean scores do not necessarily correspond to meaningful differences in user experience. Prior work emphasizes the importance of distinguishing between statistical and practical significance in IR evaluation\cite{urbano2013}. This evaluation setup approximates offline ranking evaluation commonly used in recommender systems when user interaction logs are unavailable, where ranking metrics serve as standard measures of retrieval effectiveness in the absence of user interaction data.

\paragraph{Statistical Significance Testing}
To assess whether observed differences between models are robust across queries, we perform paired t-tests on per-query NDCG@5 scores. For each proxy task, query-level scores are aligned across model pairs under a paired design. We report statistical significance at the p < 0.05 level; unadjusted p-values are reported for these exploratory pairwise comparisons, and we acknowledge that this may inflate the risk of Type I errors. This analysis allows us to determine whether observed performance differences are consistent across queries rather than arising from random variation.

% Section 4: Results
\section{Results}\label{sec:results}\label{ch:results}
On average, the Composer Identity Proxy contains 9.0 relevant tracks per query, while the Primary Musical Proxy contains 68.1 relevant tracks per query. This substantial difference in target density indicates that the musical-character proxy defines a significantly less selective retrieval task.

\subsection{Composer Identity Proxy Results}
The Composer Identity Proxy task evaluates the ability of embedding models to retrieve tracks sharing the same composer. This task serves as a controlled retrieval setting in which relevance is defined by a well-structured metadata attribute (composer identity), enabling verification that embedding models capture fundamental compositional structure. We observe consistent trends at K=10, with no qualitative changes in model ranking or statistical significance patterns, indicating that the reported findings are not specific to shallow ranking cutoffs. Notably, the relative ordering of models remains stable between K=5 and K=10, further supporting the robustness of these observations.

Table~\ref{tab:sanity_results} presents retrieval performance grouped by architectural family and ordered by capacity tier within each family.

\begin{table}[!htbp]
\centering
\caption{Composer Identity Proxy Results (Composer Retrieval). Metrics reported are corpus-level mean values at rank 5 and rank 10, including NDCG@K and Hit Rate@K (HR@K).}
\label{tab:sanity_results}
\resizebox{\columnwidth}{!}{%
\begin{tabular}{lcccc}
\hline
\textbf{Model} & \textbf{NDCG@5} & \textbf{NDCG@10} & \textbf{HR@5} & \textbf{HR@10} \\
\hline
CNN-Small & 0.548 & 0.549 & 0.640 & 0.700 \\
CNN-Medium & 0.545 & 0.555 & 0.640 & 0.709 \\
CNN-Large & 0.585$^{*\dagger}$ & 0.585$^{*\dagger}$ & 0.660$^{*\dagger}$ & 0.714$^{*\dagger}$ \\
Transformer-Medium & 0.642$^{*}$ & 0.639$^{*}$ & 0.709$^{*}$ & 0.749$^{*}$ \\
Transformer-Large & 0.588$^{\dagger}$ & 0.582$^{\dagger}$ & 0.665$^{\dagger}$ & 0.719$^{\dagger}$ \\
\hline
\end{tabular}
}
\vspace{2pt}
\raggedright
\footnotesize{
$^{*}$ denotes statistically significant difference compared to the smallest model within the same architectural family.
$^{\dagger}$ denotes statistically significant difference compared to the preceding capacity tier (paired t-test, $p < 0.05$).
}
\end{table}

The Composer Identity Proxy results suggest that model capacity can influence ranking quality under structured relevance definitions when relevance is defined by a relatively structured and objective criterion. Within the CNN family, small differences in NDCG@5 are observed across capacity tiers, with the large model achieving the highest mean performance. However, these differences are relatively modest in magnitude. Within the Transformer family, the medium model achieves higher performance than the large model in both NDCG@5 and HR@5. Some of these differences are statistically significant (paired t-test, p < 0.05), particularly when compared to CNN-based models.

On the composer proxy, Transformer-Medium achieves the highest mean NDCG@5 and outperforms all other models. These differences are statistically significant for several pairwise comparisons, particularly against CNN-based models (paired t-test, p < 0.05). This suggests that the observed improvements are consistent across queries rather than driven by a small number of examples, indicating robust ranking performance under this structured retrieval setting.

\subsection{Primary Musical Proxy Results}
The Primary Musical Proxy task evaluates retrieval based on abstract musical character labels (e.g., "Energetic" vs. "Calm"). This task represents the core retrieval challenge for the archive.

Table~\ref{tab:musical_results} summarizes the performance across all models.

\begin{table}[!htbp]
\centering
\caption{Primary Musical Proxy Results (Character Retrieval). Metrics reported are corpus-level mean values at rank 5 and rank 10, including NDCG@K and Hit Rate@K (HR@K).}
\label{tab:musical_results}
\resizebox{\columnwidth}{!}{%
\begin{tabular}{lcccc}
\hline
\textbf{Model} & \textbf{NDCG@5} & \textbf{NDCG@10} & \textbf{HR@5} & \textbf{HR@10} \\
\hline
CNN-Small & 0.652 & 0.649 & 0.783 & 0.798 \\
CNN-Medium & 0.656 & 0.656 & 0.768 & 0.798 \\
CNN-Large & 0.642 & 0.643 & 0.749 & 0.783 \\
Transformer-Medium & 0.632 & 0.617 & 0.778 & 0.793 \\
Transformer-Large & 0.631 & 0.626 & 0.763 & 0.783 \\
\hline
\end{tabular}
}
\vspace{2pt}
\raggedright
\footnotesize{
No statistically significant differences are observed between model variants (paired t-test, $p > 0.05$).
}
\end{table}

In contrast to the Composer Identity Proxy, the Musical Proxy results do not show a clear benefit from increased capacity. No statistically significant differences are observed across model variants in either the CNN or Transformer families (paired t-test, p > 0.05), 
indicating that their performance is indistinguishable under this evaluation setting. 
Although minor variations in mean scores are observed, these differences should not be interpreted as systematic improvements or declines.

These results suggest that, within this evaluation setting, embedding capacity alone does not consistently improve ranking quality when relevance is defined by proxy labels derived from abstract musical character annotations. From a recommender systems perspective, these findings indicate that increasing model capacity does not consistently improve candidate ranking quality under the evaluated conditions.

For the Primary Musical Proxy, no pairwise differences are statistically significant (paired t-test, p > 0.05 for all comparisons), indicating that observed performance differences are not consistent across queries and are likely driven by variability rather than systematic model improvements.

\subsection{Capacity--Performance Trends Within Architectural Families}

This section analyzes scaling behavior strictly within architectural families to isolate the effect of parameter count from architectural inductive bias. Taken together, the within-family scaling patterns described below provide the primary empirical test of the condition-dependent capacity hypothesis. These findings suggest that improvements in model capacity do not consistently translate into improved candidate ranking quality within the evaluated retrieval setting

\subsubsection{CNN Scaling Behavior}

Within the CNN family, capacity scaling exhibits task-dependent behavior. On the Composer Identity Proxy, performance differences are observed across capacity tiers, with the large model achieving the highest mean NDCG@5. Some of these differences are statistically significant (p < 0.05), indicating that higher-capacity models can provide measurable gains under this structured retrieval setting. On the Musical Proxy, no statistically significant differences are observed across capacity tiers (p > 0.05), indicating that performance differences are not systematic under this evaluation setting. These results suggest that the effect of capacity scaling depends on the structure of the retrieval task within the evaluated dataset.

\subsubsection{Transformer Scaling Behavior}

Within the Transformer family, the effect of scaling also depends on the proxy task. On the Composer Identity Proxy, the medium model achieves higher performance than the large model, indicating that increased capacity does not necessarily lead to improved ranking quality even under structured retrieval conditions. On the Musical Proxy, however, the medium and large transformer models perform comparably, indicating that increased capacity does not translate into improved performance when evaluated on this abstract semantic proxy. This contrast reinforces the view that transformer scaling benefits are conditional rather than universal in the present setting. These results suggest that scaling effects differ between structured retrieval criteria and abstract semantic similarity tasks.

\subsection{Cross-Family Descriptive Comparison}

Comparing performance across families reveals that scaling effects depend more on task structure than on raw model size alone. On the Composer Identity Proxy, performance differences are observed across capacity tiers in both families, although these effects are not consistently monotonic. On the Musical Proxy, however, neither family exhibits consistent improvement, and performance differences remain small.

These comparisons are descriptive and do not imply causal superiority of one architecture over another, since architecture, parameter count, and pretraining regime co-vary across families. The narrow performance range across models suggests limited representational separability under the current retrieval setting, regardless of capacity tier.

Given the small dataset size, absolute differences should be interpreted cautiously, with emphasis placed on consistent trends rather than isolated values. From a recommender systems perspective, these results indicate that architectural differences and scaling do not substantially alter candidate ranking quality within the constraints of the current dataset

\subsection{Embedding Extraction Cost}

Latency is measured as the average processing time per track in the embedding extraction pipeline. 
All latency benchmarks were conducted on an Apple Mac Mini M1 (2020) with 16GB RAM, using sequential inference on the CPU without GPU acceleration and a batch size of 1.

To complement the ranking results, we also report embedding extraction latency for each model during offline feature generation.

\begin{table}[!htbp]
\centering
\caption{Embedding extraction latency across model capacities.}
\label{tab:latency_results}
\begin{tabular}{lc}
\hline
\textbf{Model} & \textbf{Latency (ms/track)} \\
\hline
CNN-Small & 2178.87 \\
CNN-Medium & 3108.93 \\
CNN-Large & 4217.94 \\
Transformer-Medium & 23145.73 \\
Transformer-Large & 55724.24 \\
\hline
\end{tabular}
\end{table}

Table~\ref{tab:latency_results} shows that embedding extraction cost increases with model capacity, particularly for the Transformer family. We emphasize that the reported latency corresponds to the offline embedding extraction process during catalog ingestion, rather than real-time retrieval latency. Once embeddings are computed and indexed, online candidate retrieval is performed via similarity search, which is independent of the underlying model used for feature extraction. Combined with the ranking results, these results suggest that larger models may not provide a favorable trade-off between computational cost and retrieval effectiveness under the current experimental setup and hardware conditions.


\subsection{Summary of Empirical Findings}

The experimental results yield three main observations. First, increasing model capacity does not consistently improve ranking performance across tasks. On the Composer Identity Proxy, some statistically significant differences are observed between model variants, while on the Musical Proxy, no statistically significant improvements are observed across capacity tiers. Second, scaling behavior differs by task. Structured composer-based retrieval exhibits measurable differences across model capacities, whereas abstract semantic proxy tasks do not show statistically significant variation, indicating that the effect of scaling depends on the retrieval setting. Third, performance values cluster within a relatively narrow range across models, suggesting that model capacity is not the dominant factor influencing retrieval quality in this small-scale dataset.

Collectively, these findings support a condition-dependent interpretation of scaling behavior. The impact of embedding-model capacity varies depending on the structure of the retrieval task and the nature of the proxy relevance labels. In this evaluated setting, task formulation and label characteristics appear to play a more influential role than capacity alone in determining ranking quality. 

From an information retrieval perspective, these results highlight that ranking effectiveness may be influenced by the definition and reliability of proxy relevance signals, in addition to model capacity.

% Section 5: Discussion
\section{Discussion}

\subsection{RQ1: Capacity Effects in Small-Scale Retrieval}

The results suggest that, within this dataset and evaluation setting, embedding-model capacity does not consistently improve candidate ranking quality. This suggests that the benefits of scaling are highly task-dependent, particularly for abstract semantic tasks such as musical character retrieval. From a recommender systems perspective, this implies that improvements in representation capacity do not necessarily translate into improved ranking effectiveness when relevance signals are weakly structured.

This observation is consistent with prior work showing that strong MIR performance does not necessarily translate to improved recommendation quality \cite{tamm2024}. In addition, the stylistic coherence of classical music within a single archive may lead to a limited degree of representational separability, where different pieces are mapped to similar regions of the representation space regardless of model capacity. From an information retrieval perspective, this indicates that ranking effectiveness may be influenced by the reliability and structure of proxy relevance labels, in addition to model capacity. Importantly, these trends remain consistent across ranking depths (K=5 and K=10), indicating that the observed effects are not sensitive to the choice of cutoff. A key factor underlying these observations is the difference in target density between the two proxy tasks. On average, the Composer Identity Proxy contains 9.0 relevant tracks per query, whereas the Primary Musical Proxy contains 68.1 relevant tracks per query. This substantial difference in target density indicates that the musical-character proxy defines a significantly less discriminative and less selective retrieval task.

The substantially higher number of relevant items in the musical-character proxy suggests that this task is less sensitive to differences in representation quality, which may partially explain the weaker sensitivity to model capacity observed in this setting. This highlights a key limitation of metadata-derived proxy tasks: when relevance definitions are broad, improvements in representation quality may not be reflected in ranking metrics.

\subsection{RQ2: Conditions for Capacity Benefits}

The results suggest that higher-capacity models may provide measurable benefits within this setting primarily when the retrieval task is aligned with structured and well-defined attributes, such as composer identity.

In contrast, for abstract semantic tasks such as musical character retrieval, clear benefits from scaling are not observed. This suggests that the effectiveness of capacity scaling depends on the alignment between representation space and the structure of the relevance signal.

From a ranking perspective, larger models may provide improved separability when relevance is well-defined in some cases, but fail to provide additional discrimination when labels are coarse, subjective, or weakly informative.

\subsection{Hypothesis Evaluation}

The empirical results support the hypothesis that increasing embedding-model capacity does not necessarily improve candidate ranking quality in this setting. Instead, performance is primarily governed by the interaction between model capacity and the structure and reliability of proxy relevance signals rather than capacity alone. Importantly, this pattern remains consistent across ranking depths (K=5 and K=10), indicating that the observed effects are robust to the choice of evaluation cutoff. However, future work is needed to determine how dataset characteristics mediate this relationship.

\subsection{Implications for Recommender System Design}

Because our evaluation relies on metadata-derived proxy relevance labels rather than historical user interaction logs, these findings reflect constraints observed in offline retrieval evaluation settings commonly used in real-world systems when user interaction data are unavailable. They suggest that, absent collaborative signals, increasing embedding capacity may not reliably improve candidate generation in small-scale or data-constrained settings.

Higher-capacity models introduce additional computational cost without proportional gains in ranking performance. In practical systems, this implies that lightweight or mid-sized models may provide a favorable cost-performance trade-off under the evaluated conditions for candidate generation. However, it is important to note that the reported latency measurements are obtained under CPU-bound, sequential inference on a Mac Mini M1 without hardware acceleration. Transformer-based models are known to benefit substantially from parallelized execution on GPUs or TPUs, as well as from batched inference. Consequently, the observed computational cost differences are specific to CPU-bound deployment settings and may differ in production environments with hardware acceleration.

Computational cost is explicitly evaluated in Section~4.5 (Table~\ref{tab:latency_results}), where embedding extraction latency increases with model capacity, particularly for Transformer-based models. Combined with the ranking results, this suggests that higher-capacity models may not provide a favorable trade-off under the current experimental setup and hardware conditions between computational cost and retrieval effectiveness in this setting. The consistency of these findings across multiple ranking depths further supports their robustness under different retrieval evaluation settings.

More broadly, these results highlight that ranking quality may be influenced by factors such as task formulation, label quality, and dataset structure, in addition to model capacity. For archival systems such as iPalpiti, this suggests prioritizing metadata quality, semantic labeling, and dataset curation over scaling model size.

\subsection{Limitations}
This study has several limitations. First, the dataset is small and restricted to a single-domain classical music archive, which limits generalizability. Second, the musical-character labels used in the proxy task are coarse and partially subjective, which constrains the evaluation of fine-grained semantic similarity. Third, the evaluation is based on constructed proxy relevance labels rather than real user interaction data, and therefore may not fully reflect real-world recommendation effectiveness.

Finally, the pseudo-labels used in the primary proxy task are partially derived from a pretrained emotion model. Although these labels are used solely as a fixed evaluation reference and do not share training objectives with the evaluated models, this setup may still introduce bias. Unlike supervised evaluation settings where labels and model objectives are aligned, the pseudo-labels in this study are generated by an external model trained on a distinct task (music emotion recognition), while the evaluated embedding models are trained using different objectives (e.g., audio tagging or self-supervised representation learning). Therefore, the evaluation does not measure agreement with the training objective of the embedding models, but rather assesses whether their learned representations correlate with independently derived semantic descriptors. While this reduces the risk of direct methodological circularity, the reliance on automatically generated labels remains a limitation, and future work should incorporate human-annotated ground truth for more robust evaluation.

Importantly, this study does not aim to establish general scaling laws for audio representation learning, but rather to characterize behavior under a specific data-constrained, single-domain setting. The findings should therefore be interpreted as context-dependent observations rather than universally generalizable conclusions.

While statistical significance testing is performed, the observed differences between model variants remain relatively small and dependent on the proxy task, and should therefore be interpreted with caution given the limited dataset and evaluation setting.

\subsection{Directions for Future Work}
Future work includes validating these findings on larger and more diverse datasets to assess generalizability, as well as incorporating human-annotated relevance signals to reduce reliance on pseudo-labels. Additionally, exploring richer and more fine-grained semantic labels may improve the evaluation of abstract similarity in music retrieval.

Finally, integrating user interaction signals and conducting user studies will be important for understanding whether improvements in embedding representations translate into better recommendation outcomes in real-world settings.


% Section 6: Conclusion
\section{Conclusion}

This study investigated the effect of embedding-model capacity on candidate ranking quality in a content-based music retrieval setting for a small-scale classical music archive. Through controlled experiments across CNN and Transformer models, we find that, within the evaluated setting, increasing model capacity does not consistently improve candidate ranking quality. While capacity-related differences are observable for structured retrieval tasks, these effects are limited in magnitude and remain consistent across ranking depths (K=5 and K=10), and are not consistently observed for more automatically-derived affective proxy tasks.

These findings support a task-dependent view of scaling, where the effectiveness of model capacity is primarily governed by the structure and reliability of proxy relevance labels rather than model size alone. In particular, statistically significant differences are observed for the composer-based proxy, whereas no significant differences are found for the abstract musical-character proxy. In small, single-domain archives, ranking effectiveness is influenced more by task formulation and label quality than by embedding capacity alone.

From a system design perspective, deploying large-scale pretrained models introduces additional computational overhead. Our results suggest that, under the evaluated CPU-bound and unaccelerated conditions, this additional cost does not consistently translate into improved ranking effectiveness. In such resource-constrained settings, smaller or mid-sized models may provide a favorable cost-performance trade-off for offline feature extraction and catalog ingestion.

% Manual bibliography workflow:
% Edit CIKM_refs.tex directly, then compile CIKM.tex twice with pdflatex.
% Manual references for CIKM.tex
\begin{thebibliography}{20}

\bibitem{tamm2024}
Y.-M. Tamm and A. Aljanaki, ``Comparative analysis of pretrained audio representations in music recommender systems,'' in \emph{Proceedings of the 18th ACM Conference on Recommender Systems (RecSys '24)}, Bari, Italy, Oct. 2024, pp. 1088--1092. doi: 10.1145/3640457.3688172.

\bibitem{zaman2023}
K.~Zaman, M.~Sah, C.~Direkoglu, and M.~Unoki, ``A survey of audio classification using deep learning,'' \emph{IEEE Access}, vol.~11, pp. 102\,345--102\,381, 2023.

\bibitem{jiang2024}
X.~Jiang, Y.~Zhang, G.~Lin, and L.~Yu, ``Music emotion recognition based on deep learning: A review,'' \emph{IEEE Access}, vol.~12, pp. 157\,716--157\,745, 2024.

\bibitem{lin2024}
P.-C.~Lin \emph{et~al.}, ``Prototype for a personalized music recommendation system based on TL-CNN-GRU model with 10-fold cross-validation,'' 2024.

\bibitem{pourmoazemi2024}
N.~Pourmoazemi and S.~Maleki, ``A music recommender system based on compact convolutional transformers,'' \emph{Expert Systems with Applications}, vol.~255, p. 124473, 2024.

\bibitem{gong2021ast}
Y.~Gong, Y.-A. Chung, and J.~Glass, ``Ast: Audio spectrogram transformer,'' \emph{arXiv preprint arXiv:2104.01778}, 2021.

\bibitem{krichene2020}
W.~Krichene and S.~Rendle, ``On sampled metrics for item recommendation,'' in \emph{Proceedings of the 26th ACM SIGKDD Conference on Knowledge Discovery and Data Mining (KDD '20)}.\hskip 1em plus 0.5em minus 0.4em\relax Virtual Event, CA, USA: ACM, 2020, pp. 1748--1757.

\bibitem{canamares2020}
R.~Ca\~namares and P.~Castells, ``On target item sampling in offline recommender system evaluation,'' in \emph{Proceedings of the 14th ACM Conference on Recommender Systems (RecSys '20)}.\hskip 1em plus 0.5em minus 0.4em\relax Virtual Event, Brazil: ACM, 2020, pp. 259--268.

\bibitem{porcaro2022}
L.~Porcaro, E.~G\'omez, and C.~Castillo,
``Assessing the impact of music recommendation diversity on listeners: A longitudinal study,''
\emph{ACM Transactions on Intelligent Systems and Technology}, 2020.

\bibitem{lin2025}
J.~Lin, S.~Huang, and Y.~Zhang,
``Deep neural network-based music user preference modeling, accurate recommendation, and IoT-enabled personalization,''
2025.

\bibitem{abbattista2024}
D.~Abbattista, V.~W. Anelli, T.~Di~Noia, C.~Macdonald, and A.~Petrov, ``Enhancing sequential music recommendation with personalized popularity awareness,'' in \emph{Proceedings of the 18th ACM Conference on Recommender Systems (RecSys '24)}, Bari, Italy, Oct. 2024.

\bibitem{schedl2021}
M.~Schedl, C.~Bauer, and W.~Reisinger, ``Listener modeling and context-aware music recommendation based on country archetypes,'' \emph{Frontiers in Artificial Intelligence}, vol.~3, 2021.

\bibitem{eerola2011}
T.~Eerola and J.~K. Vuoskoski, ``A comparison of the discrete and dimensional models of emotion in music,'' \emph{Psychology of Music}, vol.~39, no.~1, pp. 18--49, 2011.

\bibitem{urbano2013}
J.~Urbano, M.~Schedl, and X.~Serra, ``Evaluation in music information retrieval,'' \emph{Journal of Intelligent Information Systems}, 2013.

\bibitem{kong2020panns}
Q.~Kong, Y.~Cao, T.~I. Li, Y.~Wang, and M.~D. Plumbley,
``PANNs: Large-scale pretrained audio neural networks for audio pattern recognition,''
\emph{IEEE/ACM Transactions on Audio, Speech, and Language Processing},
vol.~28, pp. 2880--2894, 2020.

\bibitem{li2023mert}
Y.~Li, X.~Chen, and others, ``MERT: Acoustic Music Understanding Model with Large-Scale Self-supervised Training,'' \emph{arXiv preprint arXiv:2306.00107}, 2023.

\bibitem{schedl2018}
M.~Schedl, H.~Zamani, C.-W.~Chen, Y.~Deldjoo, and M.~Elahi,
``Current challenges and visions in music recommender systems research,''
\emph{International Journal of Multimedia Information Retrieval},
vol.~7, pp.~95--116, 2018.
doi: 10.1007/s13735-018-0154-2.

\bibitem{deldjoo2024}
Y.~Deldjoo, M.~Schedl, and P.~Knees,
``Content-driven music recommendation: Evolution, state of the art, and challenges,''
\emph{Computer Science Review},
vol.~51, p.~100618, 2024.
doi: 10.1016/j.cosrev.2024.100618.

\bibitem{sun2019bert4rec}
F.~Sun, J.~Liu, J.~Wu, C.~Pei, X.~Lin, W.~Ou, and P.~Jiang,
``BERT4Rec: Sequential recommendation with bidirectional encoder representations from transformer,''
in \emph{Proceedings of the 28th ACM International Conference on Information and Knowledge Management (CIKM '19)},
Beijing, China, 2019, pp. 1441--1450.

\bibitem{yilmaz2008}
E.~Yilmaz and J.~A. Aslam,
``Estimating average precision with incomplete and imperfect judgments,''
in \emph{Proceedings of the 17th ACM Conference on Information and Knowledge Management (CIKM '08)},
Napa Valley, CA, USA, 2008, pp. 1021--1030.

\end{thebibliography}


\end{document}
